\subsection{AXIOMS OF QUANTIFIED THEORIES}

A quantified theory is any theory \(\mathfrak{C}\) in which the schemes S1 to S4 (§3, no. 1) and the scheme S5 below provide implicit axioms.

S5. If R is a relation in C, if T is a term in C, and if x is a letter, then the relation \((T|x)R \Longrightarrow (\exists x)R\) is an axiom.

This rule is indeed a scheme. For let A be an axiom of C obtained by applying S5; there is then a relation R in C, a term T in C, and a letter x such that A is \((T|x)R \Longrightarrow (\exists x)R\) . Let U be a term in C and let y be a letter. We shall show that \((U|y)A\) can also be obtained by applying S5. Using CS1 (§1, no. 2) and CS8 (no. 1), we can confine ourselves to the case in which x is distinct from y and does not appear in U. Let \(R'\) be the relation \((U|y)R\) and \(T'\) the term \((U|y)T\) . The criteria CS2 (§1, no. 2) and CS9 (no. 1) show that \((U|y)A\) is identical with \((T'|x)R' \Longrightarrow (\exists x)R'\) .

The scheme S5 says that if there is an object T for which the relation R, considered as expressing a property of x, is true, then R is true for the object \(\tau_{x}(R)\) ; this accords with the intuitive meaning we have attributed to \(\tau_{x}(R)\) (§1, no. 3, Remark).

